% Generated by M_Tools/analysis/build_semantic_scale_paper_figures.py
% Do not edit figure paths or captions by hand; regenerate instead.

\begin{figure}[t]
\centering
\includegraphics[width=0.95\linewidth]{generated/figures/fig_iclr95_gate_scores.pdf}
\caption{ICLR 9.5 evidence gate generated from existing audits. Green bars pass the corresponding gate; orange bars remain blocked. The figure shows that problem anatomy, literature breadth, and practicality are at the target, while method and submission gates remain open.}
\label{fig:iclr95-gate}
\end{figure}

\begin{figure}[t]
\centering
\includegraphics[width=0.95\linewidth]{generated/figures/fig_support_law_correlations.pdf}
\caption{Semantic-scale support risk is predictive across OVD and closed-set remote-sensing settings.  Bars show the dataset-level correlation between continuous Gaussian support risk and high-confidence pair error.}
\label{fig:support-law}
\end{figure}

\begin{figure}[t]
\centering
\includegraphics[width=0.95\linewidth]{generated/figures/fig_deployment_utility.pdf}
\caption{Deployment utility from existing result summaries.  The top panel shows AP deltas; the bottom panel shows total and review-queue SISE reduction rates where the denominator is nonzero.}
\label{fig:deployment-utility}
\end{figure}

\begin{figure}[t]
\centering
\includegraphics[width=0.95\linewidth]{generated/figures/fig_bass_boundary_ap_risk.pdf}
\caption{BASS-GSF boundary evidence.  The scatter plot uses only rows with real eval JSON plus deployment-risk summaries.  It shows why the current train-time delta result cannot be promoted to a 9.5 method claim: reducing risk-weighted log-z must also preserve AP-sensitive ranking.}
\label{fig:bass-boundary}
\end{figure}
